


\section{Omni-WorldSuite}
\label{sec:interworld_suite}

To enable a comprehensive analysis of the interactive response capabilities of world models, \textbf{Omni-WorldSuite} constructs targeted evaluation prompts across diverse interaction levels and scenario types. In this section, we detail the construction pipeline of Omni-WorldSuite, provide representative examples, and present its statistical analysis.

\subsection{Construction Pipeline}


% The prompts in \textbf{Omni-WorldSuite} are constructed along two primary dimensions. 
% First, to capture interaction modeling—the core capability of world models—we categorize interactions by their scope of influence, namely a single object, a local region, or the global environment, forming three interaction levels. 
% Second, to ensure applicability to both general-purpose video generation models and domain-specific world models, the prompts cover both real-world general scenes and representative task-oriented scenes, such as autonomous driving, embodied robotics, and gaming environments. 
% By default, each evaluation prompt consists of a textual prompt describing the world evolution process and an initial frame image specifying the starting world state. Fig.~\ref{fig:promptsuite-pipeline} (a) and (b) illustrate two strategies for constructing the evaluation prompts.

The prompts in \textbf{Omni-WorldSuite} are designed along two primary dimensions. The first dimension is scene coverage, spanning both general daily-life scenarios and task-oriented environments such as autonomous driving, embodied AI, and gaming. Collectively, these scenarios cover key aspects of world modeling, including physical laws, commonsense reasoning, causality, camera motion, closed-loop dynamics, and spatial constraints. The second dimension is a three-level interaction hierarchy that characterizes the scope of interaction effects (Fig.~\ref{fig:intro} (Left)). Level 1 involves actions whose effects are confined to the acting object, without altering other objects or the surrounding environment. Level 2 includes localized interactions where one object directly affects another. Level 3 captures more complex interactions that influence multiple objects and lead to broader environmental changes. Each prompt is defined by a textual description of interaction-driven world-state evolution and an initial frame image specifying the starting world state. For a subset of prompts that require explicit camera control, we additionally provide camera trajectories to constrain the viewpoint transition during generation. 
Fig.~\ref{fig:promptsuite-pipeline}(a) and (b) illustrate two prompt construction strategies.


% As shown in ~\ref{fig:promptsuite-pipeline}(a), Considering that images generated by the aforementioned strategy may not always satisfy our evaluation requirements, we further adopt the strategy illustrated in Fig.~\ref{fig:promptsuite-pipeline} (a), which constructs prompts directly from real-world datasets to validate robustness across heterogeneous environments. 
% By grounding each test case in a high-fidelity initial state derived from established open-source datasets, this pipeline bypasses arbitrary text-only prompts. 
% Specifically, our data composition includes three domains: 
% \begin{itemize}
%     \item \textit{Autonomous Driving}, utilizing sequences from DriveLM~\cite{sima2024drivelm} to extract the ego-vehicle's first-frame perspectives paired with recorded camera motion vectors to evaluate the model's ability to extrapolate road dynamics.
%     \item  \textit{Embodied Robotics}, incorporating manipulation-centric tasks from InternData-A1~\cite{cai2026internvla} to observe physical causalities resulting from robot-object contact.
%     \item \textit{Gaming and Simulation}, integrating Sekai~\cite{li2025sekai} to test the preservation of consistent motion laws in non-photorealistic, highly dynamic visual spaces. 
% \end{itemize} 
% To bridge the gap between raw pixels and semantic instructions, we isolate the first frame of a given sequence as the ground-truth anchor and employ Qwen-VL~\cite{Qwen3-VL} to generate a base description of the world evolution process. Recognizing the tendency for vision-language models to occasionally hallucinate spatial relations, we enforce a strict manual verification stage. During this process, human annotators meticulously audit and refine each caption to ensure exact correspondence with the visual scene. The final synthesized prompt—comprising the verified caption, original camera trajectory, and initial frame—serves as the benchmark's grounding input.

\begin{figure}
    \centering
    \includegraphics[width=\linewidth]{Images/prompt-suite.pdf}
    \caption{Omni-WorldSuite Construction Pipeline and Analysis. \textbf{(a)} Dataset-grounded prompt generation. Prompts are generated from open-source data using first-frame and camera-motion cues, refined through VLM captioning, and finally verified by human annotators. \textbf{(b)} Concept-driven prompt generation. Prompts are derived from interaction prototypes using LLM/VLM-based generation and human curation, together with generated or edited first frames. \textbf{(c)} Suite taxonomy across indoor scenes, including diffusion (Diff.), sliding, and building-related (Buil.) scenarios; outdoor scenes, including natural, projectile motion (Proj.), and urban scenarios (Urban); and task-oriented settings, including robotics (Robot), autonomous driving (Driv.), and gaming (Game). \textbf{(d)} Coverage comparison by prompt modality and capability axes. Abbr.: \textbf{Traj} (camera trajectory); \textbf{AD} (autonomous driving), \textbf{EAI} (embodied AI); \textbf{PP} (physical principles); \textbf{LCC} (loop-closure consistency); \textbf{Cau.} (causality); \textbf{CS} (common sense); \textbf{Inter.} (Interaction).}
    \label{fig:promptsuite-pipeline}
\end{figure}

\paragraph{Dataset-grounded Prompt Generation.} As shown in Fig.~\ref{fig:promptsuite-pipeline}(a), we introduce a dataset-grounded prompt construction strategy to address the limited realism, complexity, and robustness of synthetic images. We first extract the camera motion trajectory and the first video frame from open-source datasets to serve as the motion and visual prompts, respectively. Next, we employ Qwen-VL~\cite{Qwen3-VL} to generate an initial caption for the sequence. To mitigate potential errors in spatial relations and object attributes, all generated captions are manually verified and refined to ensure consistency with the source sequence. The final evaluation prompt consists of the verified caption, the initial frame, and, when available, the original camera trajectory, serving as the grounded input for benchmark evaluation. Specifically, Omni-WorldSuite covers three domains: 
\begin{itemize}
    \item \textit{Autonomous Driving}, which uses sequences from DriveLM~\cite{sima2024drivelm}. We extract the first-frame ego-view image together with recorded camera trajectories to evaluate the model's ability to extrapolate road dynamics under realistic driving conditions. 
    \item \textit{Embodied Robotics}, which uses manipulation-oriented tasks from InternData-A1~\cite{cai2026internvla} to examine physical causality arising from robot--object interactions. 
    \item \textit{Gaming and Simulation}, which uses Sekai~\cite{li2025sekai} to test whether the model can preserve coherent motion patterns in highly dynamic and non-photorealistic environments. 
\end{itemize}



% To systematically evaluate the world model's understanding of diverse physical laws and interaction paradigms, we establish a hierarchical taxonomy categorizing interactive scenarios into four primary domains, each defined by characteristic entities and functional action sets. 
% \begin{itemize}
%     \item  The \textit{Kitchen \& Food} domain is designed to test the model's ability to handle irreversible state changes and fluid dynamics. By centering on entities such as refrigerators, kettles, and perishable items, we evaluate actions like opening/closing, pouring liquids, and slicing, probing whether the model can maintain volume conservation and texture consistency. 
%     \item The \textit{Daily Life} domain emphasizes environmental transitions and soft-body deformations. Featuring household objects like lamps and curtains, it focuses on actions such as switching lights and drawing curtains to accurately manifest global illumination changes, complex occlusions, and the intricate folding behaviors of fabrics.
%     \item The \textit{Digital \& Office} domain examines fine-grained, high-precision interactions within electronic interfaces. Through entities such as laptops and smartphones, we assess rapid-response actions like typing, scrolling, and erasing, challenging the temporal synchronization between rapid movements and dynamic screen feedback. 
%     \item Finally, the \textit{Outdoor \& Workshop} domain introduces rigid-body mechanics and functional tool manipulation. This category spans from urban navigation involving vehicles and traffic lights to manual craftsmanship using hammers and planks, demanding an understanding of momentum, collision physics, and goal-oriented tool usage.
% \end{itemize}


\paragraph{Concept-driven Prompt Generation.} As shown in Fig.~\ref{fig:promptsuite-pipeline}(b), we introduce a concept-driven prompt construction strategy featuring a generate--verify--refine pipeline to synthesize text, first frames (representing the initial world state), and camera motion trajectories.
%
Specifically, we first \textit{build a set of prototype concepts} spanning scene domains, target objects, and actions under different interaction levels. 
%
We then randomly sample an interaction level, scene type, target entity, and action from the resulting taxonomy. Conditioned on these attributes, ChatGPT-5.2~\cite{choi2025chatgpt} \textit{generates a textual prompt and a camera trajectory}. Both outputs are further cross-checked by Gemini~\cite{google2025gemini3} and DeepSeek-R1~\cite{guo2025deepseek}, followed by careful human verification and refinement. This manual revision process eliminates linguistic ambiguity and ensures the clarity, motion plausibility, and overall consistency of the evaluation cases.
%
Finally, we adopt a \textit{multi-stage image generation} pipeline to ensure high-fidelity initial frames. We use FLUX.1-dev~\cite{flux1kreadev2025} to generate $3$ candidates per prompt with a CFG scale of $3.5$ and $50$ sampling steps. All candidates are manually screened for physical plausibility, instruction adherence, and visual quality. If no valid result is obtained, we rewrite the prompt with ChatGPT-5.2 and, when necessary, apply Qwen-Image~\cite{wu2025qwenimagetechnicalreport} for refinement or artifact correction. Only minor localized in-painting is allowed during post-processing. All final images must satisfy quality control requirements, including a minimum resolution of $1024 \times 1024$, consistency with the prompt, and clear visibility of the target interactive objects. 


\begin{figure}
    \centering
    \includegraphics[width=\linewidth]{Images/dataset-examples.pdf}
  \caption{\textbf{Omni-WorldSuite examples across three interaction levels.} \textbf{Left:} Examples from the \textit{General Scene} domain. \textbf{Right:} Examples from the \textit{Task-Oriented Scene} domain, including optics/camera trajectories, game, physics/common sense, autonomous driving, and embodied AI. Each example pairs a first-frame grounding (and trajectory, when applicable) with an action prompt, with red boxes indicating interaction-relevant entities.}
\label{fig:dataset-examples}
\end{figure}
\paragraph{Omni-WorldSuite Examples.}
As Fig.~\ref{fig:dataset-examples} illustrates, we pair initial frames with action-driven prompts to demonstrate the three-level interaction hierarchy, visually anchoring relevant entities with red boxes.
\begin{itemize}
    \item \textbf{Level 1:} Actions are confined to the acting object without altering other objects or the environment. \textit{General Scenes} evaluate phenomena like physical optics (e.g., viewing fields through a crystal ball), while \textit{Task-Oriented Scenes} test continuous spatial navigation (e.g., moving along a riverside path).
    \item \textbf{Level 2:} One object directly affects another. Examples include testing thermodynamics in \textit{General Scenes} (e.g., heating a metal rod in a campfire) and complex ego-vehicle navigation alongside dynamic traffic in \textit{Task-Oriented Scenes} (e.g., autonomous driving).
    \item \textbf{Level 3:} Actions influence multiple objects and lead to broader environmental changes. Prompts cover physical causality in \textit{General Scenes} (e.g., snapping spaghetti, tidying a room) and multi-stage manipulation in \textit{Task-Oriented Scenes} (e.g., a robotic arm grasping a bottle and handing it to a person).
\end{itemize}






\subsection{Omni-WorldSuite Analysis and Statistics}


\paragraph{Concept Set Analysis.}
As shown in Fig.~\ref{fig:promptsuite-pipeline}(c), the set of prototype concepts mainly covers two broad scene categories, namely indoor and outdoor scenes, as well as task-oriented scenarios such as autonomous driving, embodied robotics, and gaming. Within each broad category, we further include several representative interaction types. Overall, these prompts span multiple dimensions, ranging from natural environments, urban scenes, and architectural spaces to fundamental physical motion, fluid and thermal phenomena, optical effects, material deformation, commonsense reasoning, object affordance, robotic manipulation, and embodied interaction, thereby forming a comprehensive prompt set that balances scene diversity, physical realism, and task interactivity. Beyond static scene descriptions, the collection also includes a large number of dynamic processes, causally driven events, and goal-oriented manipulation tasks, enabling a systematic evaluation of a model’s capabilities in scene understanding, physical consistency, spatial constraint reasoning, and embodied task execution.

To facilitate the computation of evaluation metrics, we further provide auxiliary metadata for each prompt. (i) First, we enumerate all entity objects appearing in the prompt and categorize them into affected and unaffected sets according to the interaction actions. For affected entities, we additionally annotate the expected coarse motion direction and magnitude. (ii) Next, based on the world evolution described in the textual prompt, we extract a list of key events ordered by their temporal occurrence. (iii) Finally, to evaluate camera motion and spatial consistency, we annotate expected camera motions for a subset of prompts, including the motion direction and magnitude. We also incorporate a challenging return-to-origin setting, where the model is required to return the camera to its original position after completing a motion cycle. Video frames in which the camera revisits the same spatial position are referred to as revisit frames.


\paragraph{Compare with other Benchmarks.}
As shown in Fig.~\ref{fig:promptsuite-pipeline}(d), compared with prior benchmarks such as VBench~\cite{zheng2025vbench2}, WorldScore~\cite{duan2025worldscore}, and WorldModelBench~\cite{li2025worldmodelbench}, Omni-WorldBench supports the most comprehensive set of prompt modalities, encompassing text, image, and trajectory inputs. Moreover, it evaluates both task-oriented and general scenes, rather than focusing on only a narrow subset of scenarios. Specifically, it covers a diverse range of scene and reasoning types, including physical regularities, loop-closure motion, causal reasoning, and commonsense reasoning, thereby achieving the broadest coverage of scenario types among existing benchmarks. Furthermore, Omni-WorldBench is the first benchmark to explicitly account for interaction types as a core evaluation dimension. This comprehensive design provides a reliable testbed for the development and evaluation of next-generation 4D world models.

\begin{figure}
    \centering
    \includegraphics[width=1\linewidth]{Images/distribution.pdf}
    \caption{Statistics of Omni-WorldSuite. 
    \textbf{(a)} Overall Distributions; \textbf{(b--g)} Distributions of core principles; \textbf{(h)} prompt counts by interaction level; \textbf{(i--k)} word clouds of objects, actions, and scenes. NM (Newtonian Mechanics), FM (Fluid Mechanics), MP (Material Properties), WO (Waves and Optics), MC (Momentum and Collision), TP (Thermodynamics and Phase Transition), EC (Energy Conversion and Conservation); SEK (Scene/Event Knowledge), OFK (Object Function Knowledge), HAK (Human Action Knowledge); C2B (Condition-to-Behavior), A2M (Action-to-Motion), C2O (Collision-to-Outcome); TFS (Tracking / Follow Shot), OAS (Orbit / Arc Shot), HHS (Handheld / Shaky); ART (Axial Round-Trip Motion), ODC (Optical / Dynamic Consistency Closure), SCC (Spiral / Composite Closure), CDR (Curved / Diagonal Return Motion), PCP (Planar Closed-Path Motion), ORC (Orbital / Rotational Closure), UNC (Uncategorized); MKC (Mechanical / Kinematic Constraints), CSS (Contact \& Support Stability), OAC (Occlusion \& Accessibility Constraints), CL (Containment \& Leakage), GFSC (Geometric Fit \& Size Compatibility), DMC (Deformation \& Material Constraints).}
    \label{fig:suite-static}
\end{figure}

\paragraph{Statistics.}
Omni-WorldSuite contains 1,068 evaluation prompts, making it a comparatively large evaluation suite for video generation assessment. As shown in Fig.~\ref{fig:suite-static}(a), the suite exhibits a multi-label distribution over six major annotation dimensions, namely Physics Principles (PP), Commonsense (CS), Causality (Cau), Camera Motion (CM), Loop-Closure Consistency (LCC), and Spatial Constraints (SC). Among these dimensions, Physics Principles appears most frequently, followed by Causality and Commonsense. Fig.~\ref{fig:suite-static}(b--g) further present the subcategory distributions within each dimension. Specifically, NM and FM are the most frequent categories in Physics Principles; SEK dominates the Commonsense dimension; C2B is the most common causal type; Pan and Tilt are the most frequent camera motion patterns; ART and ODC are the most common loop-closure categories; and MKC appears most frequently among the spatial constraints. Fig.~\ref{fig:suite-static}(h) further shows that Level 2 contains the largest number of prompts, followed by Level 3 and Level 1. In addition, the word clouds in Fig.~\ref{fig:suite-static}(i--k) highlight the diversity of objects, actions, and scenes covered by the suite. Overall, these statistics indicate that Omni-WorldSuite is not only large in scale, but also broad in semantic and structural coverage, providing a diverse testbed for evaluating interactive world modeling under physical, causal, spatial, and motion-related constraints.


%0303
% 当前世界模型的应用范围已从客观物理场景扩展到具身智能、自动驾驶和游戏环境等任务导向场景。尽管应用场景多样，这些模型的共同本质都是对给定交互条件下环境状态时序演化的建模。因此，统一评测的关键不应仅在于场景类型的划分，更应关注交互过程本身的复杂度。基于这一认识，我们将 Omni-WorldSuite 按照交互影响范围与因果传播复杂度组织为三个层级，分别对应单一目标状态变化、局部环境响应以及跨实体与时间的全局状态传播，并据此构建了覆盖自然、城市等通用环境以及具身、自驾、游戏等典型任务导向场景的 1011 条 prompts。
% The application scope of current world models has expanded from objective physical scenes to task-oriented domains such as autonomous driving, embodied robotics, and gaming environments. Despite the diversity of application settings, these models share a common essence: modeling the temporal evolution of environmental states under given interaction conditions. Consequently, the key to unified evaluation should hinge not merely on categorizing scene types, but rather on capturing the complexity of the interaction process itself. Based on this insight, we organize Omni-WorldSuite into \textcolor{blue}{three hierarchical levels according to the scope of interaction effects and the complexity of causal propagation: single-target state changes, local environmental responses, and global state propagation across entities and time.} Following this taxonomy, we construct 1,011 prompts covering both general environments (e.g., natural and urban scenes) and typical task-oriented scenarios, including embodied interaction, autonomous driving, and gaming.

% Omni-WorldSuite contains \textbf{1011 prompts} spanning diverse domains (natural scenes, autonomous driving, embodied/robotics tasks, and game/simulated environments) to mitigate domain-specific bias and promote robust comparison across models.
% Prompts are stratified by the Level I/II/III taxonomy to cover self-state consistency, localized target transitions, and global cascading evolution.
% To quantify coverage, we summarize the distribution across levels and domains in \ref{tab:suite_stats} and report additional attribute statistics (e.g., number of entities involved, propagation horizon, and physical/common-sense mechanisms) in \ref{fig:suite_attr}.
% Empirically, we observe a consistent performance decay as tasks progress from Level I to Level III, supporting that the taxonomy reflects increasing causal complexity and providing a principled basis for Omni-Metric to attribute errors to (i) incorrect target transitions, (ii) non-target instability, or (iii) broken propagation chains.
% The full prompt list and metadata are released with our toolkit; we also include a complete listing in App.~\ref{app:full_list} for reproducibility.


% The evaluation prompts in \textbf{Omni-WorldSuite} are constructed along two primary and orthogonal dimensions. 
% First, to capture interaction modeling—the core capability of world models—we categorize interactions by their scope of influence, namely a single object, a local region, or the global environment, forming three interaction levels. 
% Second, to ensure applicability to both general-purpose video generation models and domain-specific world models, the prompts cover both real-world general scenes and representative task-oriented scenes, such as autonomous driving, embodied robotics, and gaming environments. 
% By default, each evaluation prompt consists of a textual prompt describing the world evolution process and an initial frame image specifying the starting world state. 
% Fig.~\ref{fig:promptsuite-pipeline} (a) and (b) illustrate two strategies for constructing the evaluation prompts.


% % The construction pipeline is as follows:
% % \begin{figure}
% %     \centering
% %     \includegraphics[width=1\linewidth]{Images/prompt-pipeline.pdf}
% %     \caption{Omni-World Suite Construction Pipeline.}
% %     \label{fig:promptsuite-pipeline}
% % \end{figure}



% For the strategy shown in Fig.~\ref{fig:promptsuite-pipeline} (a), we first build a set of \emph{prototype concepts} according to interaction levels and scene categories. 
% The construction of scene categories and their associated typical entities and actions is detailed in App.~\ref{app:prompt_schema}. 
% Based on these concepts, we adopt a \textit{generate--verify--refine} pipeline to construct the textual prompts.
% Specifically, (i) we randomly sample an interaction level, scene type, target entity, and action from the prototype concept set; 
% (ii) ChatGPT-5 then generates an initial textual prompt conditioned on these keywords; 
% (iii) to reduce linguistic ambiguity and ensure consistency with the evaluation requirements, each prompt is cross-checked by Gemini and DeepSeek-R1. The filtered prompts are further manually edited and refined to ensure clarity and executability; 
% (iv) given the initial world state implied by the textual prompt, we employ powerful image generation and editing models, including FLUX~1.0, Qwen-Image, and ChatGPT-5.2, to produce multiple candidate images, from which the most suitable image is selected through human review.
% Details of the instruction prompts used for different MLLMs, as well as the implementation specifics of the image generation and processing models, are provided in App.~\ref{app:prompt_schema} and App.~\ref{app:special_domains}.


% Omni-WorldSuite is built with a reproducible ``generate--verify--refine'' pipeline designed to maximize executability and minimize ambiguity. As shown in Fig.~\ref{fig:promptsuite-pipeline}, we first define interaction prototypes consistent with hree hierarchical levels according to the scope of interaction effects and the complexity of causal propagation: single-target state changes, local environmental responses, and global state propagation across entities and time, then use \textbf{ChatGPT-5} to draft prompts in standardized natural language.
% To reduce underspecification and linguistic ambiguity, each prompt is cross-checked by \textbf{Gemini} and \textbf{DeepSeek-R1} for consistency, missing conditions, and observability of intended state changes, followed by human editing on a per-prompt basis.
% We provide the exact prompt schema, annotation fields, and representative examples in App.~\ref{app:prompt_schema}.
% rules in App.~\ref{app:first_frame_policy,app:special_domains}.



% \paragraph{Camera-motion prompts.}
% A subset of prompts explicitly tests camera motion and spatial consistency. Beyond checking motion accuracy (e.g., pan/tilt/dolly trajectories), we include ``return-to-origin'' variants that require the model to recover the original viewpoint/composition after a motion cycle, exposing accumulated drift and long-horizon inconsistency.
% The full camera-motion taxonomy and criteria are provided in App.~\ref{app:camera_motion}.

% Considering that images generated by the aforementioned strategy may not always satisfy our evaluation requirements, we further adopt the strategy illustrated in Fig.~\ref{fig:promptsuite-pipeline} (b), which constructs prompts from real-world datasets. 
% Specifically, (i) we first curate several public datasets from environments that require high realism and structural consistency, including autonomous driving, embodied robotics, and gaming (see App.~\ref{app:first_frame_policy} for the dataset list), and select appropriate initial-frame images. 
% (ii) Based on these images, we employ Qwen-VL to generate textual prompts describing the world evolution process, which are subsequently verified and refined by human annotators to obtain the final prompts.


% To facilitate the computation of evaluation metrics, we further provide auxiliary metadata for each prompt. 
% (i) First, we enumerate all entity objects appearing in the prompt and categorize them into affected and unaffected sets according to the interaction actions. For affected entities, we additionally annotate the expected coarse motion direction and magnitude. 
% (ii) Next, based on the world evolution described in the textual prompt, we extract a list of key events ordered by their temporal occurrence. 
% (iii) Finally, to evaluate camera motion and spatial consistency, we annotate expected camera motions for a subset of prompts, including the motion direction and magnitude. 
% We also incorporate a challenging \emph{return-to-origin} setting, where the model is required to return the camera to its original position after completing a motion cycle. 
% Video frames in which the camera revisits the same spatial position are referred to as \emph{revisit frames}. 
% Implementation details and visualization examples are provided in App.~\ref{app:prompt_schema}.


